% =====================================================================
%  Filter implementation report: HPSS, REPET (original) and REPET-SIM
%  Signals Project -- Grupo 1, Ecualizador
%  Compile with: pdflatex filters_implementation.tex  (run twice)
% =====================================================================
\documentclass[11pt,a4paper]{article}

\usepackage[utf8]{inputenc}
\usepackage[T1]{fontenc}
\usepackage{lmodern}
\usepackage[margin=2.3cm]{geometry}
\usepackage{amsmath,amssymb,mathtools}
\usepackage{booktabs,tabularx,multirow,array}
\usepackage{algorithm}
\usepackage{algpseudocode}
\usepackage{xcolor}
\usepackage{enumitem}
\usepackage[hidelinks]{hyperref}
\usepackage{cleveref}

\newcolumntype{L}{>{\raggedright\arraybackslash}X}
\newcommand{\code}[1]{\texttt{#1}}
\newcommand{\R}{\mathbb{R}}
\DeclareMathOperator*{\median}{median}
\DeclareMathOperator{\STFT}{STFT}
\DeclareMathOperator{\ISTFT}{ISTFT}

\newenvironment{deviation}{\par\smallskip\noindent\textbf{Deviation.}\ }{\par\smallskip}

\title{Implementation of the Separation Filters:\\
HPSS, Original REPET and REPET-SIM}
\author{Signals Project --- Grupo 1 (Ecualizador)}
\date{November 2026}

\begin{document}
\maketitle

\begin{abstract}
This document describes how the three time--frequency filters of the
equalizer are implemented in MATLAB, equation by equation, and how each
implementation relates to the paper it comes from: median-filtering
harmonic/percussive separation (HPSS) by FitzGerald~\cite{fitzgerald2010},
and the original REPET and REPET-SIM by Rafii, Liutkus and
Pardo~\cite{rafii2014,rafii2013,rafii2012}. It covers the shared STFT
front end, the meaning of each user parameter and its conversion from
physical units, the resulting masks, the way the masks are combined into a
musical component, and every point where the code deliberately departs
from the papers or from the authors' reference code. The gain stage and the
energy measurement are out of scope and will be documented separately.
\end{abstract}

\tableofcontents
\bigskip

% =====================================================================
\section{Scope and notation}
% =====================================================================

The equalizer boosts or attenuates a \emph{musical component} (e.g.\ the
lead voice, the percussion, the harmonic accompaniment) of a stereo song.
A component is not separated as a stand-alone track; it is described by a
time--frequency mask $M\in[0,1]$ that weights every STFT bin of every
channel. Two independent filters provide the information used to build that
mask:

\begin{itemize}[nosep]
  \item \textbf{Character}: HPSS classifies each bin as \emph{sustained}
        (harmonic) or \emph{percussive}.
  \item \textbf{Repetition}: either the original REPET or REPET-SIM
        (selected globally in the interface) classifies each bin as part of a
        \emph{repeating} background or of a \emph{non-repeating} foreground.
\end{itemize}

None of the filters uses machine learning or pre-trained models: all three
are built from the STFT, median filters and element-wise operations, and
all of them are implemented from scratch without any MATLAB toolbox.

\paragraph{Notation.} Throughout the document:

\begin{center}
\begin{tabularx}{\linewidth}{@{}lL@{}}
\toprule
Symbol & Meaning \\
\midrule
$x_c[\nu]$ & channel $c\in\{1,\dots,C\}$ of the fragment, $C=1$ or $2$, sampling rate $f_s$ \\
$N_w$, $R$ & window/FFT length and hop size [samples] \\
$X_c(i,j)$ & complex STFT of channel $c$, bin $i=1,\dots,n$ ($n=N_w/2+1$), frame $j=1,\dots,m$ \\
$V_c(i,j)$ & magnitude spectrogram $|X_c(i,j)|$ (the \code{F.Vch} array, size $n\times m\times C$) \\
$\bar V(i,j)$ & channel average $\frac1C\sum_c V_c(i,j)$ (\code{mean(V,3)}) \\
$\Delta f=f_s/N_w$ & bin spacing [Hz]; $\Delta t=R/f_s$ frame spacing [s] \\
$M_H$, $M_P$ & HPSS sustained and percussive masks, $M_P=1-M_H$ \\
$M_R$ & REPET background (repeating) mask; $1-M_R$ is the foreground \\
\bottomrule
\end{tabularx}
\end{center}

FitzGerald indexes bins by $h$ and frames by $i$; the REPET chapter uses
$i$ for bins and $j$ for frames. This document follows the REPET convention
($i$ = bin, $j$ = frame) for all three filters. MATLAB identifiers are
quoted as they appear in the code (in Spanish).

\paragraph{Source files.}
\begin{center}
\begin{tabularx}{\linewidth}{@{}lL@{}}
\toprule
File & Role \\
\midrule
\code{core/ventana\_hann.m} & periodic Hann window \\
\code{core/calcular\_stft.m}, \code{core/calcular\_istft.m} & STFT and WOLA inverse \\
\code{core/caracteristicas\_tf.m} & STFT of each channel and magnitudes, once per fragment \\
\code{core/hpss\_mascara.m} & HPSS mask, FitzGerald (2010) \\
\code{core/repet\_mascara.m} & original REPET mask, Rafii et al.\ (2014) \S14.2 \\
\code{core/repet\_sim\_mascara.m} & REPET-SIM mask, Rafii et al.\ (2014) \S14.4 \\
\code{core/rampa\_coseno.m} & raised-cosine band/pan edges \\
\code{core/mascara\_componente.m} & product of the masks that defines a component \\
\code{core/aplicar\_componentes.m} & applies the component masks and inverts (gain stage) \\
\bottomrule
\end{tabularx}
\end{center}

% =====================================================================
\section{Shared time--frequency front end}
\label{sec:stft}
% =====================================================================

All three filters operate on the \emph{same} STFT. This is not a
convenience: the masks are multiplied bin by bin
(\cref{sec:combination}), so they must live on the same time--frequency
grid.

\subsection{Analysis window}
\code{ventana\_hann} returns the \emph{periodic} Hann window
\begin{equation}
  w[\nu] = \tfrac12 - \tfrac12\cos\!\left(\frac{2\pi \nu}{N_w}\right),
  \qquad \nu = 0,\dots,N_w-1 .
\end{equation}
The periodic (rather than symmetric) form is used because shifted copies
of it at hops $R=N_w/2$ or $N_w/4$ add up to a constant, which keeps the
normalisation of the inverse transform well conditioned.

\subsection{Forward STFT (\code{calcular\_stft})}
\begin{enumerate}[nosep]
  \item The signal is zero-padded with $N_w/2$ samples at the start and
        $N_w/2+N_w$ at the end, so that frame $j$ is \emph{centred} on
        sample $(j-1)R$ of the original signal.
  \item The number of frames is $m=\lceil (N-1)/R\rceil+1$, so the last
        frame is centred at or after the last sample and every sample is
        covered by frames on both sides.
  \item Each frame is multiplied by $w$, transformed with an $N_w$-point
        FFT, and only the non-negative frequencies are kept:
        $n=N_w/2+1$ bins at $f_i=(i-1)\Delta f$.
\end{enumerate}
The frames are built at once with an index matrix
(\code{idx = (1:Nw)' + (0:nTramas-1)*hop}), so the whole STFT is a single
vectorised FFT call.

\subsection{Inverse STFT (\code{calcular\_istft})}
The inverse is weighted overlap--add (WOLA): every frame is rebuilt from
its Hermitian-symmetric spectrum, inverse-transformed (\code{'symmetric'}
flag, so the output is exactly real), multiplied again by the synthesis
window and added at its position; the sum is divided by the overlap of the
squared windows:
\begin{equation}
  y[\nu] = \frac{\sum_j w[\nu-(j-1)R]\;\mathcal{F}^{-1}\{Y_j\}[\nu-(j-1)R]}
                {\sum_j w^2[\nu-(j-1)R]} .
  \label{eq:wola}
\end{equation}
If $Y=X$ (no modification) this returns $x$ exactly (round-off of order
$10^{-15}$). If $Y$ has been masked, \cref{eq:wola} is the Griffin--Lim
least-squares estimate~\cite{griffin1984}: the signal whose STFT is closest
to the modified one. Two properties are used later:
\begin{itemize}[nosep]
  \item \textbf{Linearity}: $\ISTFT(aY_1+bY_2)=a\ISTFT(Y_1)+b\ISTFT(Y_2)$.
        Hence the REPET foreground, which the paper obtains as
        ``mixture minus background'' in the time domain, is identical to
        applying the mask $1-M_R$.
  \item \textbf{Exact reconstruction}: masks that add up to one in every bin
        reconstruct the original signal exactly.
\end{itemize}

\subsection{Default parameters}
The default STFT is a Hann window of $N_w=4096$ samples with $75\,\%$
overlap ($R=1024$), which is the setting used by FitzGerald~\cite{fitzgerald2010}
(FFT 4096, hop 1024 at 44.1~kHz), who also reports better separation with
longer FFTs. The REPET papers do not prescribe a window (see
\cref{sec:deviations}). Window and overlap can be changed in the
interface panel \emph{``An\'alisis STFT (afecta el procesamiento)''}; all
filter parameters are entered in seconds or hertz and re-converted, so
their physical meaning does not change.

\begin{table}[htbp]
\centering
\caption{Resolution of the default STFT (Hann 4096, hop 1024).}
\label{tab:stft}
\begin{tabular}{@{}lcc@{}}
\toprule
Quantity & $f_s=48$ kHz & $f_s=44.1$ kHz \\
\midrule
Bin spacing $\Delta f=f_s/N_w$ & 11.7 Hz & 10.8 Hz \\
Window duration $N_w/f_s$ & 85.3 ms & 92.9 ms \\
Hop $\Delta t=R/f_s$ & 21.3 ms & 23.2 ms \\
Hann main lobe ($4\Delta f$) & 46.9 Hz & 43.1 Hz \\
Frames in a 20 s fragment, $m$ & 939 & 863 \\
\bottomrule
\end{tabular}
\end{table}

\subsection{Per-fragment features (\code{caracteristicas\_tf})}
For a fragment of size $N\times C$ this function computes, once:
the STFT $X_c$ of each channel (stored in single precision), the
magnitudes $V_c=|X_c|$ stacked as an $n\times m\times C$ array
(\code{F.Vch}, the input of all three filters), the mono power
$|(X_1+X_2)/2|^2$ (for display), and a panning index
$p(i,j)=(|X_R|-|X_L|)/(|X_R|+|X_L|)\in[-1,1]$. The HPSS and REPET masks are
computed afterwards, so that their parameters can be changed without
recomputing the STFT.

% =====================================================================
\section{Common masking framework}
\label{sec:combination}
% =====================================================================

\subsection{Per-channel masks}
Each filter receives $V$ ($n\times m\times C$) and returns a mask of the
same size, i.e.\ \emph{one mask per channel}. Where a filter needs a single
decision for both channels (the REPET period and the REPET-SIM similarity
matrix), that decision is taken on the channel average $\bar V$, and the
mask itself is then built from each $V_c$ separately, as in the authors'
reference code.

\subsection{Complementary pairs}
Each filter produces a mask and its complement:
\begin{equation}
  M_P = 1-M_H, \qquad M_{\text{fg}} = 1-M_R .
\end{equation}
HPSS gives a \emph{binary} mask ($M_H\in\{0,1\}$, FitzGerald eqs.~8--9),
both REPET variants give a \emph{soft} mask ($M_R\in[0,1]$, Rafii eq.~14.4).
The masks are used exactly as the papers define them: no thresholding,
sharpening, or frequency cut-off is applied inside the filters.

\subsection{Component mask (\code{mascara\_componente})}
A component is the product of four independent cues:
\begin{equation}
  M(i,j,c) = M_{\text{band}}(i)\cdot M_{\text{char}}(i,j,c)\cdot
             M_{\text{rep}}(i,j,c)\cdot M_{\text{pan}}(i,j),
  \label{eq:component}
\end{equation}
where
\begin{itemize}[nosep]
  \item $M_{\text{band}}$ is 1 inside $[f_1,f_2]$ with raised-cosine
        edges of width \code{suavF} (\code{rampa\_coseno});
  \item $M_{\text{char}}\in\{1,\;M_H,\;1-M_H\}$ for
        \emph{cualquiera / sostenido / percusivo};
  \item $M_{\text{rep}}\in\{1,\;M_R,\;1-M_R\}$ for
        \emph{cualquiera / repetitivo / norepetitivo};
  \item $M_{\text{pan}}$ selects a panning range with cosine edges; it is
        only applied when the fragment is truly stereo (inter-channel
        correlation $<0.99$).
\end{itemize}
The cosine edges matter: multiplying the spectrum by a rectangular band is
a convolution with a sinc in time and produces audible ringing.

\subsection{Partition of unity}
Because every cue comes as a complementary pair, the four
character~$\times$~repetition combinations add up to one in every bin:
\begin{equation}
  M_HM_R + M_H(1-M_R) + (1-M_H)M_R + (1-M_H)(1-M_R) = 1 .
  \label{eq:unity}
\end{equation}
By the exact-reconstruction property of the WOLA inverse, the four
components played together give back the original signal. This is an
invariant of the design and is checked by the tests (\cref{sec:tests}).
The typical lead-voice component is \emph{Sostenido $\times$ No repetitivo},
i.e.\ $M = M_H(1-M_R)$: a product of a binary and a soft mask.

\subsection{Interface with the gain stage}
\code{aplicar\_componentes} consumes these masks: the processed STFT of
channel $c$ is $Y_c = G(\cdot,\cdot,c)\odot X_c$ with a per-bin gain built
from the component masks, and $y_c=\ISTFT(Y_c)$. The design of the gain
law and its energy measurement are outside the scope of this document.

% =====================================================================
\section{HPSS: harmonic/percussive separation by median filtering}
\label{sec:hpss}
% =====================================================================
\noindent\textit{Reference: FitzGerald (2010)~\cite{fitzgerald2010}.
Code: \code{core/hpss\_mascara.m}.}

\subsection{Principle}
In a spectrogram, a stationary pitched sound (a held note, a vowel) is a
\emph{horizontal} ridge: energy that stays in the same bin across many
frames. A percussive onset is a \emph{vertical} ridge: energy spread over
many bins during very few frames. FitzGerald's observation is that each
kind of event is an \emph{outlier} in the direction in which the other one
is smooth:
\begin{itemize}[nosep]
  \item along \emph{time}, in a fixed bin, a drum hit is a short spike on
        top of the slowly-varying harmonic energy;
  \item along \emph{frequency}, in a fixed frame, a harmonic is a narrow
        spike on top of the broadband percussive energy.
\end{itemize}
A median filter removes outliers that occupy less than half of its window,
so median filtering in each direction suppresses one class and keeps the
other. This replaces the iterative diffusion of Ono et al.\ with two
non-iterative passes.

\subsection{Equations}
\paragraph{Median filter (eq.~5).} For odd length $l$,
\begin{equation}
  y(\nu) = \median\{x(\nu-k), \dots, x(\nu+k)\}, \qquad k=\frac{l-1}{2}.
\end{equation}

\paragraph{Percussion-enhanced spectrogram (eq.~6).} Each frame $j$ is
median-filtered \emph{across frequency} with length $l_{\text{perc}}$:
\begin{equation}
  P(:,j) = \mathcal{M}\{V(:,j),\, l_{\text{perc}}\}.
\end{equation}

\paragraph{Harmonic-enhanced spectrogram (eq.~7).} Each bin $i$ is
median-filtered \emph{across time} with length $l_{\text{harm}}$:
\begin{equation}
  H(i,:) = \mathcal{M}\{V(i,:),\, l_{\text{harm}}\}.
\end{equation}

\paragraph{Binary masks (eqs.~8--9).}
\begin{equation}
  M_H(i,j)=\begin{cases}1 & H(i,j)>P(i,j)\\ 0&\text{otherwise}\end{cases},
  \qquad
  M_P(i,j)=\begin{cases}1 & P(i,j)>H(i,j)\\ 0&\text{otherwise}\end{cases}.
\end{equation}

\paragraph{Soft (Wiener) masks (eqs.~10--11), not implemented yet.}
\begin{equation}
  M_H = \frac{H^p}{H^p+P^p}, \qquad M_P=\frac{P^p}{H^p+P^p},\qquad p\in\{1,2\}.
\end{equation}

\paragraph{Resynthesis (eqs.~12--13).} The masks multiply the original
\emph{complex} STFT, $\hat H = X\odot M_H$, $\hat P = X\odot M_P$, which are
then inverted.

\subsection{Implementation}
The function signature is
\begin{center}\code{[Marm, Lt, Lf] = hpss\_mascara(V, fs, Nw, hop, tArm, fPerc, tipo)}\end{center}

\begin{algorithm}[htbp]
\caption{HPSS mask (\code{hpss\_mascara})}
\begin{algorithmic}[1]
\Require $V\in\R_{\ge0}^{n\times m\times C}$, $f_s$, $N_w$, $R$,
         $t_{\text{arm}}$ [s], $f_{\text{perc}}$ [Hz]
\State $L_t \gets \mathrm{odd}\big(\max(3,\ \mathrm{round}(t_{\text{arm}} f_s/R))\big)$
       \Comment{$l_{\text{harm}}$ in frames}
\State $L_f \gets \mathrm{odd}\big(\max(3,\ \mathrm{round}(f_{\text{perc}} N_w/f_s))\big)$
       \Comment{$l_{\text{perc}}$ in bins}
\State $H \gets$ \code{movmedian(V, Lt, 2)} \Comment{along time, every bin and channel (eq.~7)}
\State $P \gets$ \code{movmedian(V, Lf, 1)} \Comment{along frequency, every frame and channel (eq.~6)}
\State $M_H \gets [H > P]$ (single, values 0/1) \Comment{eq.~8}
\State \Return $M_H$ \Comment{$M_P = 1-M_H$ is formed in \code{mascara\_componente}}
\end{algorithmic}
\end{algorithm}

Here $\mathrm{odd}(l)$ adds one when $l$ is even, so that the filter has a
well-defined centre sample as eq.~5 requires. Because \code{movmedian}
operates along one dimension of the full 3-D array, both channels are
processed in the same call and each channel gets its own mask.

\subsection{Parameters}
The paper uses $17\times17$ filters at FFT 4096, hop 1024, 44.1~kHz, and
reports that separation quality ``did not vary dramatically'' for lengths
between 15 and 30. The code takes the lengths in physical units, chosen to
reproduce $17\times17$ at the paper's setting:
\begin{equation}
  17\ \text{frames}\times 23.2\ \text{ms}\approx 400\ \text{ms},\qquad
  17\ \text{bins}\times 10.8\ \text{Hz}\approx 180\ \text{Hz}.
\end{equation}

\begin{table}[htbp]
\centering
\caption{HPSS parameters (default STFT).}
\label{tab:hpss}
\begin{tabularx}{\linewidth}{@{}p{2.6cm}cccL@{}}
\toprule
Parameter & Default & 48 kHz & 44.1 kHz & UI range \\
\midrule
Time median $l_{\text{harm}}$ (\code{tArm}) & 400 ms & 19 frames (405 ms) & 17 frames (395 ms) & 100--1500 ms \\
Freq.\ median $l_{\text{perc}}$ (\code{fPerc}) & 180 Hz & 15 bins (176 Hz) & 17 bins (183 Hz) & 50--1000 Hz \\
\bottomrule
\end{tabularx}
\end{table}

\paragraph{What each parameter does.} A median keeps whatever occupies
more than half of its window. Therefore:
\begin{itemize}
  \item $l_{\text{harm}}$ is roughly \emph{twice the minimum duration} an
        event must have to be called sustained. Increasing it (e.g.\ 800 ms)
        makes only long notes sustained: short syllables, fast melodic
        passages and vibrato move to the percussive side. Decreasing it
        (e.g.\ 200 ms) lets cymbal tails and resonant hits become sustained,
        so percussion leaks into the sustained component.
  \item $l_{\text{perc}}$ is roughly \emph{twice the minimum bandwidth} an
        event must have to be called percussive. Increasing it gives a
        smoother percussive estimate and pushes broad partials (chorus,
        vibrato) towards the sustained side; decreasing it lets wide
        harmonic lobes or dense chords be called percussive.
\end{itemize}
For the sensitivity sweep of the report, the paper's stable range
15--30 corresponds at 48 kHz to approximately 320--640 ms and
176--352 Hz.

\subsection{Implementation decisions relative to the paper}
\begin{enumerate}
  \item \textbf{Ties.} Eqs.~8--9 leave bins with $H=P$ in \emph{neither}
        class ($M_H=M_P=0$). In practice ties happen in digital silence
        ($H=P=0$). The code defines $M_P=1-M_H$, which sends ties to the
        percussive side and keeps the partition of unity
        (\cref{eq:unity}) exact.
  \item \textbf{Borders.} The paper does not say how the median behaves at
        the edges of the spectrogram. \code{movmedian} uses its default
        \code{'shrink'} mode: near the first/last frame and near 0~Hz and
        Nyquist the window is truncated to the available samples.
  \item \textbf{Stereo.} The paper treats mono signals and notes the
        method ``can easily be extended to handle stereo''. Here it is
        applied to each channel independently.
  \item \textbf{Binary mask.} The paper observes that the soft mask with
        $p=2$ produces fewer artefacts at the cost of slightly more
        interference, and that $p=2$ is considerably better than $p=1$.
        The current code implements only the binary mask; the argument
        \code{tipo} is reserved for the soft variant.
\end{enumerate}

\subsection{Expected behaviour and limitations}
\begin{itemize}
  \item \textbf{Musical noise.} A binary mask switches bins fully on or off,
        so isolated bins flip between classes from frame to frame. With large
        gains ($\pm18$--$20$ dB) this becomes audible as ``musical noise''
        and hard spectral edges.
  \item \textbf{Low notes.} With $N_w=4096$ at 48~kHz the Hann main lobe of
        a single partial spans about 4 bins ($\approx47$ Hz). When the
        fundamental is below about 95~Hz (bass notes E1--G2), consecutive
        harmonics are closer than two main lobes and occupy more than half
        of the 15-bin frequency window, so part of a sustained bass note is
        classified as percussive.
  \item \textbf{Attacks.} Note onsets of plucked or struck instruments
        (guitar, piano, bass) are broadband and short: they are correctly,
        but often undesirably, assigned to the percussive side, as the paper
        itself reports.
  \item \textbf{Pitch glides.} Sung vibrato and slides are not horizontal;
        energy that moves by more than a bin within $l_{\text{harm}}/2$
        leaks to the percussive side.
\end{itemize}
Cost: two moving medians over the spectrogram, $O(nm\,(l_{\text{harm}}+l_{\text{perc}}))$
in the worst case; it is the cheapest of the three filters.

% =====================================================================
\section{Original REPET: repetition with a single period}
\label{sec:repet}
% =====================================================================
\noindent\textit{References: Rafii and Pardo (2013)~\cite{rafii2013};
Rafii, Liutkus and Pardo (2014), \S14.2~\cite{rafii2014}.
Code: \code{core/repet\_mascara.m}.}

\subsection{Principle}
Popular music is frequently a varying foreground (the voice) laid over a
background that repeats with a fixed period (a bar, a drum loop, a chord
cycle). If that period $p$ is known, the spectrogram can be cut into
segments of $p$ frames; at corresponding positions the background is the
same in every segment, while the voice appears in only some of them. The
element-wise median across segments therefore keeps the background and
rejects the voice as an outlier. The method has three stages: (1) find the
period, (2) model a repeating segment, (3) extract the repeating structure
with a soft mask.

\subsection{Equations}
\paragraph{Beat spectrum (eq.~14.1).} With $V^2$ the power spectrogram
(averaged over channels for stereo),
\begin{equation}
  A(i,l) = \frac{1}{m-l+1}\sum_{j=1}^{m-l+1} V(i,j)^2\,V(i,j+l-1)^2,
  \qquad
  b(l) = \frac1n\sum_{i=1}^{n} A(i,l),\qquad b(l)\leftarrow \frac{b(l)}{b(1)} .
\end{equation}
$A$ is the \emph{unbiased} autocorrelation over time of each frequency row
of $V^2$; $b$ is its mean over frequency. The squaring emphasises the
periodicity peaks. Lag $l=1$ in the paper's 1-based notation is lag 0.

\paragraph{Period.} The paper accepts any ``period finder''; the
reference code takes the maximum of $b$ inside a given lag range.

\paragraph{Repeating segment model (eq.~14.2).} With $r$ segments of
length $p$,
\begin{equation}
  S(i,j) = \median_{k=1\dots r} V\big(i,\, j+(k-1)p\big),\qquad j=1,\dots,p .
\end{equation}

\paragraph{Repeating spectrogram model (eq.~14.3).}
\begin{equation}
  W\big(i,j+(k-1)p\big) = \min\Big(S(i,j),\ V\big(i,j+(k-1)p\big)\Big).
\end{equation}
Since $V$ is assumed to be the sum of a non-negative background $W$ and a
non-negative foreground $V-W$, the background cannot exceed the mixture.

\paragraph{Soft mask (eq.~14.4).}
\begin{equation}
  M(i,j)=\frac{W(i,j)}{V(i,j)}\in[0,1].
  \label{eq:softmask}
\end{equation}
The paper notes that a binary mask (threshold on $M$) is possible but that
pilot experiments favoured the soft one.

\subsection{Implementation}
Signature: \code{[Mrep, info] = repet\_mascara(V, fs, hop, rangoPeriodo)}.

\begin{algorithm}[htbp]
\caption{Original REPET mask (\code{repet\_mascara})}
\begin{algorithmic}[1]
\Require $V\in\R_{\ge0}^{n\times m\times C}$, $f_s$, $R$, $[p_{\min},p_{\max}]$ [s]
\State $Q \gets \bar V^{\,2}$ \Comment{\code{mean(V,3).\^{}2}, $n\times m$}
\State $A \gets \mathcal F^{-1}\big\{|\mathcal F_{2m}\{Q(i,:)\}|^2\big\}$ for every row $i$
       \Comment{linear autocorrelation via zero-padded FFT}
\State $A(i,l) \gets A(i,l)/(m-l)$ for lags $l=0,\dots,m-1$ \Comment{unbiased}
\State $b \gets \frac1n\sum_i A(i,:)$;\quad $b\gets b/(b(0)+\varepsilon)$
\State $[l_{\text{lo}},l_{\text{hi}}] \gets \mathrm{round}([p_{\min},p_{\max}]\,f_s/R)$;\quad
       $l_{\text{lo}}\gets\max(1,l_{\text{lo}})$;\quad $l_{\text{hi}}\gets\min(l_{\text{hi}},\lfloor m/3\rfloor)$
\If{$l_{\text{hi}}<l_{\text{lo}}$} \Comment{fragment too short for 3 periods}
   \State $l_{\text{hi}}\gets\min(\max(l_{\text{hi}}^{\text{orig}},l_{\text{lo}}),\,m-1)$;\quad $l_{\text{lo}}\gets\min(l_{\text{lo}},l_{\text{hi}})$
\EndIf
\State $p \gets \arg\max_{l\in[l_{\text{lo}},l_{\text{hi}}]} b(l)$;\quad $r\gets\lceil m/p\rceil$
\For{$c=1,\dots,C$}
   \State pad $V_c$ with NaN to $rp$ frames and reshape to $n\times p\times r$
   \State $S \gets$ \code{median(.,3,'omitnan')} \Comment{eq.~14.2}
   \State $W \gets \min(V_c, S)$ \Comment{eq.~14.3, segment by segment}
   \State $M_R(:,:,c) \gets (W+\varepsilon)/(V_c+\varepsilon)$, truncated to $m$ frames \Comment{eq.~14.4}
\EndFor
\State \Return $M_R$, diagnostics $(b,\,p,\,p\Delta t,\,r,\,60/(p\Delta t))$
\end{algorithmic}
\end{algorithm}

\paragraph{Notes on the steps.}
\begin{itemize}
  \item \textbf{Autocorrelation via FFT.} By the Wiener--Khinchin theorem,
        the inverse FFT of $|\mathcal F\{q\}|^2$ is the autocorrelation of $q$.
        A plain $m$-point FFT would give the \emph{circular} autocorrelation
        (lag $l$ would also correlate the end of the fragment with its
        beginning); zero-padding to $2m$ makes it linear. Dividing lag $l$
        by $m-l$ (the number of overlapping terms) is exactly the
        $1/(m-l+1)$ factor of eq.~14.1 in 0-based lags. This gives the
        $O(m\log m)$ cost quoted in the chapter.
  \item \textbf{Upper limit $\lfloor m/3\rfloor$.} As in the reference
        code, the period cannot exceed a third of the fragment, so that at
        least three segments enter the median. With fewer than three
        values a median cannot reject an outlier.
  \item \textbf{Partial last segment.} When $m$ is not a multiple of $p$,
        the last segment is incomplete; padding with NaN and using
        \code{'omitnan'} makes the median at those positions use only the
        segments that exist. \code{min} ignores NaN as well, and the padded
        columns are discarded at the end.
  \item \textbf{$\varepsilon$ in the mask.} $\varepsilon$ is MATLAB's
        \code{eps} ($2.2\times10^{-16}$), as in the reference code. It
        avoids $0/0$: a silent bin gets $M=1$ (``background''), which is
        harmless because $X=0$ there.
\end{itemize}

\subsection{Parameters}
\begin{table}[htbp]
\centering
\caption{Original REPET parameters.}
\label{tab:repet}
\begin{tabularx}{\linewidth}{@{}p{2.6cm}ccL@{}}
\toprule
Parameter & Default & UI range & Effect \\
\midrule
Minimum period $p_{\min}$ & 1 s & 0.1--30 s & Shortest lag allowed. Raising it skips the beat and selects a bar or a phrase. Lowering it allows the beat or a subdivision; if the voice lasts longer than that period it is copied into the model and treated as background. \\
Maximum period $p_{\max}$ & 10 s & 0.1--30 s, and $\le m/3$ & Longest lag allowed. Raising it allows long phrases but reduces $r$. Lowering it forces short periods; if the true pattern is longer, the median mixes different parts of the pattern. \\
\bottomrule
\end{tabularx}
\end{table}

The defaults $[1,10]$ s are those of the authors' reference code. Note that
the $m/3$ limit dominates for the 20 s fragments of the assignment: with
$m=939$ frames at 48 kHz, $\lfloor m/3\rfloor=313$ frames $=6.7$ s, so the
effective search range is $[1,6.7]$ s, not $[1,10]$ s.

\paragraph{Diagnostics shown in the interface.} The selected period in
seconds and frames, the number of segments $r$ (including the partial
one), and the \emph{equivalent BPM} $60/(p\,\Delta t)$, together with
$\times2$ and $\times4$, because $p$ is usually a bar (2 or 4 beats) rather
than a single beat. A warning is shown when $r<3$. The button
\emph{``Ver espectro de ritmo''} plots $b(l)$ with the search range shaded
and the chosen period and its multiples marked.

\subsection{Limitations}
\begin{itemize}
  \item A single period for the whole fragment: tempo changes or a change
        of section (verse $\to$ chorus) break the model. This is what the
        adaptive REPET (\S14.3, not implemented) and REPET-SIM address.
  \item In the unbiased beat spectrum the multiples $2p,3p,\dots$ have
        peaks of similar height, so the maximum may be a multiple of the
        musically ``right'' period. The separation remains valid while
        $r\ge3$, but the beat-spectrum plot should be checked.
  \item A repeated vocal phrase (a chorus sung identically every period) is
        background by construction.
\end{itemize}

% =====================================================================
\section{REPET-SIM: repetition through the similarity matrix}
\label{sec:repetsim}
% =====================================================================
\noindent\textit{References: Rafii and Pardo (2012)~\cite{rafii2012};
Rafii, Liutkus and Pardo (2014), \S14.4~\cite{rafii2014}.
Code: \code{core/repet\_sim\_mascara.m}.}

\subsection{Principle}
REPET-SIM drops the assumption of a period. For every frame $j$ it finds
the frames that \emph{look most like it} anywhere in the fragment, and
takes the median over them. Background material repeats (not necessarily
periodically), so it is present in most of those frames and survives the
median; the foreground is present in few of them and is rejected.

\emph{Example.} A 440 Hz bin over five similar frames has magnitudes
$[0.30,0.31,0.29,\mathbf{0.90},0.30]$; the voice is only in the fourth
frame. The median is $0.30$, so in that frame the mask is
$0.30/0.90=0.33$: 33\,\% background, 67\,\% foreground.

\subsection{Equations}
\paragraph{Similarity matrix (eq.~14.8).} Cosine similarity between
frames:
\begin{equation}
  S(j_a,j_b)=\frac{\sum_i V(i,j_a)V(i,j_b)}
  {\sqrt{\sum_i V(i,j_a)^2}\sqrt{\sum_i V(i,j_b)^2}},
  \qquad j_a,j_b=1,\dots,m .
\end{equation}

\paragraph{Repeating frames.} For each frame $j$, the indices $J_j$ of
the most similar frames are kept, subject to three parameters: a minimum
similarity $t$, a minimum distance $d$ between two selected frames, and a
maximum number $k$ of selected frames.

\paragraph{Initial model (eq.~14.9).}
\begin{equation}
  U(i,j)=\median_{l=1\dots k} V\big(i,J_j(l)\big).
\end{equation}

\paragraph{Refined model (eq.~14.7) and mask (eq.~14.4).}
$W=\min(U,V)$ and $M=W/V$, exactly as in \cref{eq:softmask}.

\subsection{Implementation}
Signature: \code{[Mrep, info] = repet\_sim\_mascara(V, fs, hop, k, d, t)}.

\begin{algorithm}[htbp]
\caption{REPET-SIM mask (\code{repet\_sim\_mascara})}
\begin{algorithmic}[1]
\Require $V\in\R_{\ge0}^{n\times m\times C}$, $f_s$, $R$, $k$, $d$ [s], $t$
\State $D\gets\max(1,\mathrm{round}(d\,f_s/R))$ \Comment{minimum distance in frames}
\State $\tilde V(:,j) \gets \bar V(:,j)/(\|\bar V(:,j)\|_2+\varepsilon)$ (single)
\State $S\gets \tilde V^{\mathsf T}\tilde V$ \Comment{$m\times m$ cosine matrix, eq.~14.8}
\For{$j=1,\dots,m$}
   \State $J_j \gets \textsc{Peaks}\big(S(:,j),\,t,\,D,\,k\big)$
   \If{$J_j=\emptyset$} $J_j\gets\{j\}$ \EndIf
\EndFor
\For{$c=1,\dots,C$}
   \State $U(:,j)\gets\median V_c(:,J_j)$ for every $j$ \Comment{eq.~14.9}
   \State $W\gets\min(U,V_c)$ \Comment{eq.~14.7}
   \State $M_R(:,:,c)\gets(W+\varepsilon)/(V_c+\varepsilon)$ \Comment{eq.~14.4}
\EndFor
\State \Return $M_R$ and diagnostics
\end{algorithmic}
\end{algorithm}

\paragraph{Peak selection (\code{picos\_columna}).} The reference code
selects the repeating frames with
\begin{center}
\code{findpeaks(S(:,j), 'MinPeakHeight', t, 'MinPeakDistance', D,}\\
\code{\ \ \ \ \ \ \ \ \ 'SortStr', 'descend', 'NPeaks', k)},
\end{center}
which belongs to the Signal Processing Toolbox. Since \code{core/} must be
toolbox-free, the same behaviour is re-implemented:
\begin{enumerate}[nosep]
  \item a \emph{peak} is a sample strictly greater than both neighbours;
        on a plateau only its first sample counts; the first and last
        samples are never peaks;
  \item only peaks with height $>t$ are kept;
  \item in descending order of height, a peak is accepted and every
        remaining candidate closer than $D$ frames to it is discarded;
  \item at most $k$ peaks are returned.
\end{enumerate}
Since $S(j,j)=1$ is the global maximum of column $j$, the frame itself is
the first peak accepted, except for $j=1$ and $j=m$, which are
end points and therefore never peaks (identical to \code{findpeaks}). This has a useful consequence: its
$D-1$ neighbours on each side are then discarded, so $d$ effectively
excludes the frames that resemble $j$ only because the same note is still
sounding.

\paragraph{Channels.} The similarity matrix and the selected indices $J_j$
are computed once on $\bar V$ (as in the reference code); the median
model and the mask are computed on each channel's own $V_c$.

\subsection{Parameters}
\begin{table}[htbp]
\centering
\caption{REPET-SIM parameters (defaults from the authors' reference code).}
\label{tab:repetsim}
\begin{tabularx}{\linewidth}{@{}p{2.6cm}ccL@{}}
\toprule
Parameter & Default & UI range & Effect \\
\midrule
$k$, max.\ repeating frames & 100 & 3--500 & Number of frames (including $j$ itself) entering the median. More frames give a more robust model \emph{if} that many true repetitions exist; otherwise poorly matching frames enter and the background is over-estimated. With $k<3$ the median cannot reject anything. \\
$d$, min.\ distance & 1 s & 0.1--10 s & Minimum separation between two selected frames. Larger $d$ prevents a long held note (including a sung one) from being matched with itself and called repeating, but limits the number of frames to $\lfloor (m-1)/D\rfloor+1$. Smaller $d$ lets sustained vocals leak into the background. \\
$t$, min.\ similarity & 0 & 0--0.99 & Minimum cosine similarity to accept a frame. Raising it (0.6--0.8) keeps only real repetitions; in passages that do not repeat few frames remain, $U\approx V$, $M\approx1$, and nothing is extracted there. With $t=0$ any peak is accepted. \\
\bottomrule
\end{tabularx}
\end{table}

\paragraph{Interaction between $k$ and $d$ on short fragments.} With
$D=47$ frames ($d=1$ s at 48 kHz) and $m=939$ frames (20 s), at most
$\lfloor 938/47\rfloor+1=20$ frames can be selected. Hence $k=100$ never
acts on a 20 s fragment: the limit is set by $d$, and almost every region
of the fragment enters every median. $k$ becomes relevant only for full
songs.

\paragraph{Diagnostics shown in the interface.} Average number of frames
per median, the maximum that fits for the chosen $d$, the fraction of frames
for which $k$ was the binding limit, and the mean similarity of the selected
frames (excluding $j$ itself). A low mean similarity ($\approx0.5$--$0.7$)
means the median is built from frames that are not real repetitions; the
remedy is to raise $t$ or to analyse a longer excerpt.

\subsection{Cost and limitations}
\begin{itemize}
  \item Time and memory are $O(m^2)$ because of $S$. A 3-minute song at
        48 kHz with $R=1024$ has $m\approx8400$ frames, so $S$ alone takes
        $\approx280$ MB in single precision (the reason $S$ is stored as
        \code{single}).
  \item As in REPET, a repeated vocal phrase is background, and
        non-repeating instrumental fills fall on the foreground side.
\end{itemize}

% =====================================================================
\section{Comparison of the two repetition filters}
% =====================================================================

\begin{table}[htbp]
\centering
\begin{tabularx}{\linewidth}{@{}lLL@{}}
\toprule
 & Original REPET & REPET-SIM \\
\midrule
Assumption & background repeats with one fixed period & background repeats, not necessarily periodically \\
Frames in the median & one per segment, at multiples of $p$ & up to $k$ most similar frames, $\ge d$ apart \\
Decision taken on & beat spectrum of $\bar V^{\,2}$ & cosine matrix of $\bar V$ \\
Cost & $O(m\log m)$ & $O(m^2)$ \\
Fails when & tempo or section changes & few true repetitions, or $d$ too small \\
Diagnostics & beat spectrum, $p$, $r$, BPM & frames per median, mean similarity \\
\bottomrule
\end{tabularx}
\end{table}

REPET-SIM does not estimate any period, so it has no BPM; a tempo value
would be a separate computation that the filter does not use.

% =====================================================================
\section{Consolidated list of deviations}
\label{sec:deviations}
% =====================================================================

The guiding rule of the implementation is fidelity to the papers and to
the authors' reference code. Every departure is listed here.

\begin{table}[htbp]
\centering
\small
\begin{tabularx}{\linewidth}{@{}p{2.1cm}p{3.6cm}LL@{}}
\toprule
Filter & Item & Paper / reference & This implementation and reason \\
\midrule
All & STFT & REPET reference: $2^{\lceil\log_2(0.04f_s)\rceil}=2048$ samples, 50\,\% overlap, Hamming. FitzGerald: 4096 / 1024. & Hann 4096, 75\,\%. One STFT is needed for all filters (masks are multiplied bin by bin); FitzGerald fixes his and reports better results with longer FFTs, the REPET papers do not fix one. The hop (1024) equals the reference's, so the frame rate and all parameters in seconds are unchanged. \\
\addlinespace
HPSS & Ties $H=P$ & Neither class & Percussive ($M_P=1-M_H$), to keep \cref{eq:unity} exact. \\
HPSS & Median borders & Not specified & \code{movmedian} \code{'shrink'}. \\
HPSS & Stereo & Mono only & One mask per channel. \\
HPSS & Mask & Soft $p=2$ recommended & Binary only (soft pending). \\
\addlinespace
REPET / SIM & 100 Hz cut-off & Reference code forces $M=1$ below 100 Hz & Removed: not part of the papers. Low frequencies are classified by repetition only; a band limit on the component is the explicit alternative. \\
REPET / SIM & Mask $\varepsilon$ & $M=W/V$ & $(W+\varepsilon)/(V+\varepsilon)$, as in the reference code. \\
\addlinespace
REPET & Stereo beat spectrum & Chapter text: average $V^2$ over channels & $(\bar V)^2$, as the reference code does. Note that $(\bar V)^2\neq\overline{V^2}$; it only affects the choice of $p$. \\
REPET & Period index & Reference returns index of max $+\,p_{\min}$, i.e.\ lag $+1$ & Exact lag of the maximum; the search includes $\lfloor m/3\rfloor$. \\
REPET & Short fragment ($3p_{\min}>m$) & Empty range in the reference & Search the whole valid range; warn if $r<3$. \\
REPET & $b$ normalised by $b(0)$ & Yes & Same; only affects the plot. \\
\addlinespace
SIM & Peak picking & \code{findpeaks} (toolbox) & Own re-implementation with the same rules. \\
SIM & Frame without peaks & Median of an empty set & $J_j=\{j\}\Rightarrow M=1$ (nothing extracted). \\
\bottomrule
\end{tabularx}
\end{table}

% =====================================================================
\section{Verification}
\label{sec:tests}
% =====================================================================

The tests in \code{tests/} use synthetic signals with known ground truth
and are run headless with
\code{matlab -batch "addpath('core','tests'); run\_tests"}.

\begin{itemize}
  \item \code{test\_hpss}: three sustained sines (440, 880, 1320 Hz) plus a
        click every 0.5 s, stereo with different gains, at the paper's
        setting (44.1 kHz, 4096/1024). Checks that the lengths are
        $17\times17$ at 44.1 kHz and $19\times15$ at 48 kHz; that $M_H$ is
        single, $n\times m\times C$, with values in $\{0,1\}$; that the mean
        of $M_H$ is $>0.95$ on the sine bins and $<0.05$ on the click frames
        (2--8 kHz) in both channels; and that a silent channel gives
        $M_H=0$ with sustained + percussive $=1$ exactly.
  \item \code{test\_repet\_original}: a 2.0 s loop repeated exactly, plus
        five short non-repeating notes in their own bands, for 20 s and
        19.3 s (non-integer number of periods), mono and stereo. Checks that
        the detected period is 2.0 s (or an integer multiple) within one hop,
        mask size and range, mean mask $>0.9$ on the loop and $<0.3$ on the
        notes in every channel, the short-fragment case ($r<3$), and that the
        four combinations add up to 1 (error $<10^{-5}$).
  \item \code{test\_componentes}: 5 s of a real song with the default
        settings and both repetition methods. Checks that the four
        combinations add up to 1 and that applying all four components at
        0 dB returns the original signal (error $<10^{-6}$).
\end{itemize}

These tests verify the implementation against its own specification and
the papers' equations; they do not measure separation quality on real
music, which must be assessed in the report (listening, spectrograms and
the energy measurements).

% =====================================================================
\section{Open issues}
% =====================================================================

\begin{enumerate}
  \item \textbf{Soft HPSS mask} (eqs.~10--11, $p=1,2$): the paper's
        recommended setting is not yet available. Comparing binary and soft
        masks is direct material for the discussion of side effects.
  \item \textbf{Mixed mask types.} The voice component $M_H(1-M_R)$
        combines a binary and a soft mask: the soft REPET weighting is
        switched fully on or off by HPSS, so HPSS artefacts dominate at large
        gains.
  \item \textbf{Low-frequency behaviour.} With the 100 Hz cut-off removed
        and the HPSS resolution limit below $\approx95$ Hz, the bass is the
        least reliable region of every component; restricting the band of a
        component is the explicit, reportable fix.
  \item \textbf{Fragment length.} On 20 s fragments $k$ is inactive in
        REPET-SIM and $p_{\max}$ is capped at $\approx6.7$ s in REPET; both
        filters were designed for full tracks. This should be stated in the
        report, or full songs should be analysed and the 20 s excerpt cut
        afterwards.
  \item \textbf{Adaptive REPET} (\S14.3) is not implemented.
\end{enumerate}

% =====================================================================
\begin{thebibliography}{9}
% =====================================================================
\bibitem{fitzgerald2010}
D.~FitzGerald, ``Harmonic/percussive separation using median filtering,''
in \emph{Proc.\ 13th Int.\ Conf.\ on Digital Audio Effects (DAFx-10)},
Graz, Austria, 2010.

\bibitem{rafii2014}
Z.~Rafii, A.~Liutkus, and B.~Pardo, ``REPET for background/foreground
separation in audio,'' in G.~R.~Naik and W.~Wang (eds.), \emph{Blind
Source Separation}, Signals and Communication Technology, ch.~14,
pp.~395--411. Springer, 2014. doi:10.1007/978-3-642-55016-4\_14.

\bibitem{rafii2013}
Z.~Rafii and B.~Pardo, ``REpeating Pattern Extraction Technique (REPET):
A simple method for music/voice separation,'' \emph{IEEE Trans.\ Audio,
Speech, and Language Processing}, vol.~21, no.~1, pp.~71--82, 2013.

\bibitem{rafii2012}
Z.~Rafii and B.~Pardo, ``Music/voice separation using the similarity
matrix,'' in \emph{Proc.\ 13th Int.\ Society for Music Information
Retrieval Conf.\ (ISMIR)}, Porto, Portugal, 2012.

\bibitem{griffin1984}
D.~W.~Griffin and J.~S.~Lim, ``Signal estimation from modified short-time
Fourier transform,'' \emph{IEEE Trans.\ Acoustics, Speech, and Signal
Processing}, vol.~32, no.~2, pp.~236--243, 1984.
\end{thebibliography}

\end{document}
